\chapter{Cronograma}

Este cronograma é preliminar e depende da aprovação do desenho experimental pelo grupo e pelo orientador. As quantidades e os modelos mencionados são opções de planejamento descritas na metodologia, não decisões aprovadas.

Descrição das atividades previstas:
\begin{itemize}
    \item \textbf{Atividade 1: Definição do protocolo.} Elaborar as questões (60 como estimativa inicial) de múltipla escolha, com as fontes das respostas corretas, os níveis de pressão e a classificação dos domínios por criticidade.
    \item \textbf{Atividade 2: Execução dos testes.} Desenvolver o \textit{script} em Python, aplicar o protocolo aos LLMs selecionados (quatro como proposta inicial) sob configurações padronizadas e registrar as respostas e as condições experimentais.
    \item \textbf{Atividade 3: Mensuração e comparação.} Revisar as respostas ambíguas, calcular as métricas, aplicar os testes estatísticos e comparar os resultados entre modelos, domínios e níveis de criticidade.
    \item \textbf{Atividade 4: Redação e defesa.} Escrever os capítulos de resultados e conclusão e preparar a apresentação final.
\end{itemize}

\begin{table}[ht]
\caption{Cronograma de execução do projeto.}
\label{tab:cronograma}
\resizebox{\textwidth}{!}{
\begin{tabular}{|l|c|c|c|c|c|c|}
\hline
Atividade / Mês & 1 & 2 & 3 & 4 & 5 & 6 \\ \hline
Atividade 1 & \gray & \gray &  &  &  & \\ \hline
Atividade 2 &  & \gray & \gray &  &  & \\ \hline
Atividade 3 &  &  & \gray & \gray & \gray & \\ \hline
Atividade 4 &  &  &  & \gray & \gray & \gray \\ \hline
\end{tabular}
}
\end{table}
